% ID: SEL-D0048-OB-E02
% ESTADO: validada. S049.
\textbf{a)} La rama racional sólo se aplica para $x\le1$, donde su denominador nunca se anula. Cada rama es suave en su intervalo abierto. En $x=1$, $f(1)=1$ y los dos límites valen $1$, de modo que $f$ es continua en todo $\mathbb R$. Las derivadas de las ramas son $(2-x)^{-2}$ y $2x-6$. Dada la continuidad en el punto de unión, sus límites en $1$ dan las derivadas laterales $1$ y $-4$. Por tanto, $f$ es derivable en $\mathbb R\setminus\{1\}$.

\textbf{b)} En $x=0$ usamos la rama racional: $f(0)=1/2$ y $f'(0)=1/4$. La recta tangente pasa por $(0,1/2)$ con pendiente $1/4$, luego $\boxed{y=x/4+1/2}$.
